\chapter{Erasure correctness: the forward simulation}
\label{chap:correctness}

This chapter states and proves the central result of the development, the forward simulation
\code{erases\_correct}, called T5 in the repository's own numbering. It is the Lean counterpart of the
theorem MetaRocq proves in [S, Sect.~7.4] --- of which the paper's displayed
\code{simple\_erases\_correct} (p.~8:62) is the unoptimised shadow: a closed, well-typed source term
that evaluates to a value is simulated by its erasure, which evaluates in \lbox{} to an erasure of that
value. Three structural differences shape everything that follows. First, the source semantics is not
Rocq's \code{eval} but \SEval{} (Chapter~\ref{chap:source}), a model of Lean's compiler front end: it
reads a compiler-body table for its delta steps, its constructor-value arm is arity-bounded, and its
delta, iota and projection arms each carry an explicit \code{StepDefeq} certificate that the step is a
definitional equality of the kernel environment. Second, the target side is a \emph{composite}: the
erasure relation \Erases{} produces a specification \lbox{} term $t_0$, and a separate pass relation
\Lower{} (Chapter~\ref{chap:lowering}) compiles eliminator applications into \code{case} nodes and
mutual blocks into \code{fix} nodes, producing the emitted term $t$. MetaRocq has no analogue of this
second layer, and most of the work in the arms below is inverting it. Third, erasure here is a
\emph{relation} and not a function --- Lean's relevance oracle cannot be modelled functionally at this
level --- so the theorem is universally quantified over the erasure image, and the functional character
is recovered only on the first-order fragment, in the last technical section.

The nine modules the chapter owns are \code{ErasesCorrect/Steps.lean} (the induction motive, the three
step interfaces, the two statement abbreviations), \code{ErasesCorrect.lean} (the aggregator and the
eight non-step arms), \code{ErasesCorrect/Iota.lean}, \code{ErasesCorrect/Delta.lean} and
\code{ErasesCorrect/Proj.lean} (one hard arm each), \code{ErasesCorrect/Close.lean} (the theorem
itself), \code{SubjectReduction.lean}, \code{FixMetatheory.lean} (the de Bruijn convention shared by the
block-closing loop and the \lbox{} \code{fix} node) and \code{FirstOrderInd.lean} (the first-order
fragment and the observable), all under \code{LeanToLambdaBox/}. None of them contains a \code{sorry},
an \code{unsafe} or a \code{partial} declaration.

The proof is one induction on the \SEval{} derivation with the motive \code{Simulates}, which carries a
fourth conjunct --- the environment relation \ErasesEnv{}, again at the image of the \emph{value} --- as
an accumulator; without it no induction hypothesis could be applied at a contractum, since \ErasesEnv{}
is stated at a program and a contractum is a different program. Of the eleven \SEval{} rules, four are
trivial, four are discharged in the aggregator (\code{zeta}, \code{indVal}, \code{ctorVal} and the
hardest of them, \code{beta}), and three are factored out as hypothesis interfaces so that each may live
in its own module. Every arm opens the same way: the erasure inversion offers the box reading first, and
that reading is closed once and for all by \code{erases\_correct\_boxLow} and its spine form, which is
MetaRocq's box case plus Letouzey's rule $\Box\,x \to \Box$ plus the fact that $\Box$ lowers only to
$\Box$. The last technical section turns the simulation into an \emph{observation}, following
[S, Sect.~7.3]: on values of first-order inductive types the non-determinism of the relation collapses,
and the unique image is a box-free constructor tree.

\section{The statement: the motive, the step interfaces, T5}
\label{sec:correctness:statement}

\begin{definition}[The induction motive]
\label{def:Simulates}
\lean{LeanToLambdaBox.Simulates}
\leanok
\uses{def:Erases, def:ErasesEnv, def:Lower, def:WcbvEval, def:eraseFlags, def:lean4lean-grounding}
Let $e$ and $v$ be source terms. $\mathsf{Simulates}\ env\ bo\ Us\ \Gamma_{spec}\ \Gamma\ e\ v$ holds
when, for \emph{every} composite \lbox{} image of $e$ that the specification environment answers, the
value $v$ has an image the emitted program reaches: for all $ve$, $t_0$, $t$ with $e$ translating to
$ve$, \Erases{} relating $e$ to $t_0$, \Lower{} relating $t_0$ to $t$ at $\Gamma_{spec}$ and
\ErasesEnv{} holding of $t_0$, there are $v_0, v'$ with \Erases{} relating $v$ to $v_0$, \Lower{}
relating $v_0$ to $v'$, \WcbvEval{} taking $t$ to $v'$ at $\Gamma$ under \code{eraseFlags}, and
\ErasesEnv{} holding of $v_0$. The Lean definition is \code{Simulates}; its last conjunct is the
accumulator that lets the beta, zeta, delta and iota arms apply their induction hypothesis at a
contractum, and it fixes the local context at the empty one, so the motive speaks only of closed terms.
\end{definition}

\begin{definition}[The three step interfaces]
\label{def:StepIota}
\lean{LeanToLambdaBox.StepIota, LeanToLambdaBox.StepProj, LeanToLambdaBox.StepDelta}
\leanok
\uses{def:Simulates, def:SEval, def:UpstreamAsks, def:LowerEnv, def:StepDefeq, def:CasesOnShape,
def:CtorOf, def:IndArity, def:InformativeInd, def:ConstOrigin}
Each interface restates one \SEval{} rule as a hypothesis bundle: every premise of the rule, each
subderivation paired with its induction hypothesis (\code{Simulates} at that subterm), concluding
\code{Simulates} at the redex --- and each opens with \code{UpstreamAsks\ env}, the bundle of lean4lean
facts the pin does not prove. \code{StepIota} is stated at a \code{casesOn} spine
$\mkApps{}\ (\mathtt{const}\ con\ us)\ (pre \mathbin{+\!\!+} disc :: minors \mathbin{+\!\!+} extra)$
whose discriminant reaches a constructor value, the over-application $extra$ carried explicitly;
\code{StepProj} at $\mathtt{proj}\ S\ i\ disc$; \code{StepDelta} at a spine headed by a constant with a
\emph{compiler} body, unfolded at the call site's universe levels, with the side condition that the
constant is no \code{casesOn} name. All three are \code{abbrev}s, \ie{} propositions, and all three are
inhabited by theorems of this chapter (\ref{thm:step_iota}, \ref{thm:step_proj}, \ref{thm:step_delta}).
\end{definition}

\begin{remark}
The factoring is a module-graph device, not a weakening: the arm modules import the aggregator's own
dependencies, and Lean rejects the cycle a single file would need. \code{Close.lean} discharges all
three, so T5's premise list does not mention them.
\end{remark}

\begin{definition}[T5, the erasure-correctness statement]
\label{def:ErasesCorrectStmt}
\lean{LeanToLambdaBox.ErasesCorrectStmt}
\leanok
\uses{def:Simulates, def:Erases, def:ErasesEnv, def:Lower, def:LowerEnv, def:SEval, def:WcbvEval,
def:eraseFlags, def:UpstreamAsks, def:SEvalFlags, def:lean4lean-grounding}
\code{ErasesCorrectStmt} is the proposition: for every closed source term $e$ that translates to $ve$
and evaluates to $v$ under \SEval{} at source flags $fl$, and every \lbox{} program $t$ obtained by
erasing $e$ to $t_0$ and lowering $t_0$ at a specification environment $\Gamma_{spec}$ that satisfies
\ErasesEnv{} at $t_0$ and whose pruning to the emitted $\Gamma$ satisfies \LowerEnv{}, granted the
kernel environment is well formed and granted \code{UpstreamAsks\ env}, there are $v_0$ and $v'$ with
\Erases{} relating $v$ to $v_0$, \Lower{} relating $v_0$ to $v'$, and \WcbvEval{} taking $t$ to $v'$ at
\code{eraseFlags}. Eight binders carry seven premises: MetaRocq's five (\code{env.WF}, the source
typing, the source evaluation, the erasure, and \ErasesEnv{} in the role of \code{erases\_deps}), plus
\LowerEnv{} for the pass layer MetaRocq has no analogue of, plus \code{UpstreamAsks\ env}; the erasure
and the lowering are two binders so that the statement can name the middle term $t_0$ at which the
environment premise is stated. The \emph{target} flags are fixed at \code{eraseFlags}; the \emph{source}
flag record stays a parameter.
\end{definition}

\begin{definition}[T5, folded through the composite relation]
\label{def:ErasesCorrectLBStmt}
\lean{LeanToLambdaBox.ErasesCorrectLBStmt}
\leanok
\uses{def:ErasesCorrectStmt, def:ErasesLB, def:WcbvEval, def:LowerEnv, def:UpstreamAsks}
The same simulation read through the composite relation \code{ErasesLB} (source term to emitted \lbox{}
term via \emph{some} intermediate erasure), with the environment premise universally quantified over
whichever middle term the composite exhibits, and the conclusion
$\exists v'.\ \code{ErasesLB}\ v\ v' \wedge \WcbvEval{}\ t\ v'$. The Lean abbreviation is
\code{ErasesCorrectLBStmt}; this is the relation the shipping eraser's output is a member of, so it is
the form Chapter~\ref{chap:refinement} consumes.
\end{definition}

\begin{remark}
T5 as \emph{published} is weaker than what its proof establishes: the inner induction proves
\code{Simulates}, whose fourth conjunct is \ErasesEnv{} at $v_0$, and the aggregator discards that
conjunct when it exits, so a consumer wanting the environment relation at the value must re-derive it.
The module doc-comments also disagree on how to count the premises (``seven binders carrying six
premises \dots{} plus the eighth'', ``eight binders and seven premises'', ``eight premises''); all three
resolve to the eight binders above.
\end{remark}

\section{Subject reduction for the source evaluator}
\label{sec:correctness:subject}

\begin{lemma}[Spine parts translate]
\label{lem:trExprS_spine_head}
\lean{LeanToLambdaBox.trExprS_spine_head, LeanToLambdaBox.trExprS_spine_mem}
\leanok
\uses{def:lean4lean-grounding}
If an application spine $\mkApps{}\ hd\ as$ has a lean4lean translation, then its head translates, and
so does each argument separately. The Lean statements are \code{trExprS\_spine\_head} and
\code{trExprS\_spine\_mem}; both leave the local context general.
\begin{proof}
\leanok
\uses{def:lean4lean-grounding}
Induction on the argument list, peeling one application node of the translation per step; the second
form is the first applied at each prefix. These are the workhorses by which every arm obtains the typing
witness its induction hypothesis needs at an argument position.
\end{proof}
\end{lemma}

\begin{lemma}[Subject reduction along a spine, at an abstract step relation]
\label{lem:SEval-defeq_spine}
\lean{LeanToLambdaBox.SEval.defeq_spine, LeanToLambdaBox.SEval.defeq_spine_value}
\leanok
\uses{def:lean4lean-grounding, def:SEval}
Let $P$ be any relation on source terms such that whenever $e$ translates to $ev$ and $P\ e\ v$ holds,
$v$ translates to some $vv$ definitionally equal to $ev$. If a head translating to $hve$ is
definitionally equal to a second head's translation, and every argument is related by $P$ to its value,
then the whole spine's translation is definitionally equal to the value spine's. The Lean statement is
\code{SEval.defeq\_spine}; \code{SEval.defeq\_spine\_value} is its instance at a constant-headed spine
that keeps its head, which the two value arms share.
\begin{proof}
\leanok
\uses{lem:trExprS_spine_head, def:lean4lean-grounding, asm:lean4lean-trust}
Induction on the argument list. At each step lean4lean's congruence for application combines the head's
definitional equality with the argument's, and uniqueness of translations reconciles the two
translations of the same term. Leaving $P$ abstract makes the schema serve both the constructor-value
arms and the delta arm, whose side condition is stated at the \emph{evaluated} arguments.
\end{proof}
\end{lemma}

\begin{theorem}[Subject reduction as definitional equality]
\label{thm:SEval-defeq}
\lean{LeanToLambdaBox.SEval.defeq}
\leanok
\uses{def:SEval, def:lean4lean-grounding, def:StepDefeq}
If $e$ translates to $ve$ in the kernel environment and $e$ evaluates to $v$ under \SEval{}, then $v$
translates to some $vv$ with $ve$ and $vv$ definitionally equal in $env$. The Lean statement is
\code{SEval.defeq}; it is stated at a general local context and puts no restriction on the source flags.
\begin{proof}
\leanok
\uses{lem:SEval-defeq_spine, def:StepDefeq, lem:trExprS_spine_head, asm:lean4lean-trust}
Induction on the evaluation derivation. Beta is closed by lean4lean's instantiation lemma, its beta rule
for definitional equality and uniqueness of types; zeta by the corresponding instantiation for
\code{let}, at which the model side does not move at all; sort and Pi are reflexive; the two value arms
are Lemma~\ref{lem:SEval-defeq_spine}. The design point is what the \emph{other three} arms do: delta on
a compiler body, iota at a \code{casesOn} constant and projection are not kernel reductions of $env$, so
they do not prove their definitional equality --- the \SEval{} rule carries it as its \code{StepDefeq}
premise, and this theorem is where that premise is spent.
\end{proof}
\end{theorem}

\begin{remark}
The measured axiom footprint of \code{SEval.defeq} contains \code{sorryAx}, inherited from the pinned
lean4lean; it is the archetype of the trust class the repository calls B. Carrying each compiler step's
own definitional equality as a premise is a deliberate weakening relative to a Rocq-style development,
and is what lets a compiler-oriented evaluator be related to the kernel at all. The theorem is the
gating fact that lets the simulation replace a term by its value without leaving the typed world.
\end{remark}

\section{Erasure along a spine, and the box arms}
\label{sec:correctness:box}

\begin{lemma}[Erasure is a spine congruence]
\label{lem:Erases-mkApps_forall2}
\lean{LeanToLambdaBox.Erases.mkApps_forall₂}
\leanok
\uses{def:Erases, def:LBTerm-mkApps}
If the head $hd$ erases to $th$ and the arguments $as$ erase pointwise to $ts$, then
$\mkApps{}\ hd\ as$ erases to the \lbox{} spine $\mathsf{mkApps}\ th\ ts$. The Lean statement is
\code{Erases.mkApps\_forall}$_2$, whose pointwise premise is a \code{List.Forall2} witness.
\begin{proof}
\leanok
\uses{def:Erases}
Induction on the pointwise witness, applying the application rule of \Erases{} once per argument. This
is the forward direction the constructor-value, delta and iota arms use to build the image of a
contractum.
\end{proof}
\end{lemma}

\begin{lemma}[Inversion of erasure along a spine]
\label{lem:erases_mkApps_inv}
\lean{LeanToLambdaBox.erases_mkApps_inv}
\leanok
\uses{def:Erases, def:Erasable, def:lean4lean-grounding, def:LBTerm-mkApps}
If $\mkApps{}\ hd\ as$ erases to $t$ then either (a) the spine erased structurally --- head to some
$th$, arguments pointwise to $ts$, and $t = \mathsf{mkApps}\ th\ ts$ --- or (b) some \emph{prefix} of
the spine, head included, was judged irrelevant: $as$ splits as $pre \mathbin{+\!\!+} suf$, the term
$\mkApps{}\ hd\ pre$ translates to an \Erasable{} model term, the remaining arguments erase pointwise to
$ts$, and $t = \mathsf{mkApps}\ \Box\ ts$. The Lean statement is \code{erases\_mkApps\_inv}.
\begin{proof}
\leanok
\uses{def:Erases, lem:Erases-sort_inv}
Induction on the argument list; at each application node the inversion of \Erases{} offers either the
box rule or the congruence, and the box rule extends the boxed prefix by one argument. This is the exact
analogue of the two ways MetaRocq's \code{erases} can meet an application spine, and it is the case
split every spine-shaped arm opens with.
\end{proof}
\end{lemma}

\begin{lemma}[An irrelevant head makes the spine irrelevant]
\label{lem:erasable_mkApps}
\lean{LeanToLambdaBox.erasable_mkApps}
\leanok
\uses{def:Erasable, def:lean4lean-grounding}
In a well-formed environment and local context, if $\mkApps{}\ hd\ as$ translates to $ve$, the head
translates to $hve$, and $hve$ is \Erasable{}, then $ve$ is \Erasable{}. The Lean statement is
\code{erasable\_mkApps}.
\begin{proof}
\leanok
\uses{lem:trExprS_spine_head, thm:Erasable-app, asm:lean4lean-trust}
The application rule of \Erasable{} iterated along the spine, uniqueness of translations reconciling the
model terms at each node. This turns a boxed \emph{prefix} into irrelevance of the whole spine --- how
the box disjunct of Lemma~\ref{lem:erases_mkApps_inv} gets closed --- and it is also the tool the iota
and projection arms use to \emph{refute} a boxed reading of a constructor value.
\end{proof}
\end{lemma}

\begin{lemma}[The box arm, at the erasure]
\label{lem:erases_correct_box}
\lean{LeanToLambdaBox.erases_correct_box, LeanToLambdaBox.erases_correct_boxSpine}
\leanok
\uses{def:Erases, def:Erasable, def:WcbvEval, def:eraseFlags, def:SEval}
If $e$ translates, some translation of $e$ is \Erasable{}, and $e$ evaluates to $v$, then $v$ erases to
$\Box$, and $\Box$ --- already a \lbox{} value --- evaluates to an erasure of $v$. The spine form
\code{erases\_correct\_boxSpine} says the same at $\mkApps{}\ hd\ (pre \mathbin{+\!\!+} suf)$ with the
irrelevance at the prefix: the target folds $\Box\,t_1\dots t_n$ to $\Box$ by the iterated rule
$\Box\,x \to \Box$, provided each discarded argument evaluates, which the Lean statement carries as an
explicit hypothesis.
\begin{proof}
\leanok
\uses{thm:SEval-defeq, lem:erasable_mkApps, lem:WcbvEval-mkApps_box, def:Erases, asm:lean4lean-trust}
Theorem~\ref{thm:SEval-defeq} moves the irrelevance from the term to its value along the recorded
definitional equality, irrelevance being stable under definitional equality, so the box rule of
\Erases{} applies at $v$. On the target, $\Box$ evaluates to itself, and at a spine the folding lemma
for boxed applications collapses the application of $\Box$ to the evaluating arguments.
\end{proof}
\end{lemma}

\begin{lemma}[The box arm, at the composite]
\label{lem:erases_correct_boxLow}
\lean{LeanToLambdaBox.erases_correct_boxLow, LeanToLambdaBox.erases_correct_boxSpineLow}
\leanok
\uses{def:Lower, def:ErasesEnv, def:ErasesBox, def:Erases, def:WcbvEval, def:eraseFlags}
This is the form each arm actually needs, the simulation's conclusion being stated on the composite. If
$e$'s erasure is $\Box$ (an \code{ErasesBox} witness), $\Box$ lowers to $t$, and \ErasesEnv{} holds of
$\Box$, then for $e$'s value $v$ there are $v_0 = \Box$ and $v' = \Box$ with \Erases{}, \Lower{},
\WcbvEval{} taking $t$ to $\Box$, and \ErasesEnv{} at $v_0$. The spine form
\code{erases\_correct\_boxSpineLow} is the same with middle term $\mathsf{mkApps}\ \Box\ ts$ and a
hypothesis that every \emph{lowered} argument evaluates.
\begin{proof}
\leanok
\uses{lem:erases_correct_box, lem:ErasesEnv-ofReach, prop:Lower-source_const, thm:eval_deterministic,
lem:Lower-source_app, lem:WcbvEval-mkApps_box, asm:lean4lean-trust}
$\Box$ lowers only to $\Box$, so the target term \emph{is} $\Box$; Lemma~\ref{lem:erases_correct_box}
supplies an erasure of the value, and determinism of \lbox{} evaluation forces it to be $\Box$; the box
closure of \ErasesEnv{} re-establishes the accumulator, since $\Box$ reaches no kername. The spine form
first inverts the pass at the $\Box$-headed spine --- the head is neither an application nor a constant,
so no eliminator key can fire --- and then folds by the boxed-application lemma.
\end{proof}
\end{lemma}

\begin{remark}
This pair is the closing move of the box disjunct in \emph{every} arm of the simulation, and the one
place where the two-layer target shows up as pure overhead relative to [S], whose box case is one line.
\end{remark}

\section{Inverting the pass at an application}
\label{sec:correctness:lowerinv}

\begin{lemma}[The two readings of a lowered application]
\label{lem:Lower-source_app}
\lean{LeanToLambdaBox.Lower.source_app, LeanToLambdaBox.Lower.source_mkApps}
\leanok
\uses{def:Lower, def:ElimDecl, def:LBTerm-mkApps}
If \Lower{} relates $\mathtt{app}\ f\ a$ to $t$ then either (a) the pass acted as an application
congruence, $t = \mathtt{app}\ f'\ a'$ with $f$ lowering to $f'$ and $a$ to $a'$, or (b) the source term
is a \emph{saturated eliminator spine}: there are a runtime key $kn$ with \code{ElimDecl} data
$(iid, np, dp, nfs)$, a parameter block of length $dp$, a discriminant and one minor premise per
constructor, with $\mathtt{app}\ f\ a = \mathsf{mkApps}\ (\mathtt{const}\ kn)\ (pre \mathbin{+\!\!+}
disc :: minors)$, which the pass turns into a \code{case} node. At a spine whose head is neither an
application nor a constant, the companion \code{Lower.source\_mkApps} leaves only the congruence
reading, with lengths preserved and the arguments related pointwise.
\begin{proof}
\leanok
\uses{lem:Lower-ne_fix_of_block, def:Lower}
Case analysis on the \Lower{} derivation. The eliminator arm is indexed by an application spine, so an
eliminator node with a non-empty over-application re-associates into the congruence reading; the arm
emitting a block's \code{fix} node from a lambda body is excluded because an application is not a
lambda, and the arm emitting it from a constant because an application is not a constant. The spine form
iterates the first along the spine, the eliminator disjunct being unavailable at the excluded heads.
\end{proof}
\end{lemma}

\begin{remark}
Disjunct (b) is exactly what Proposition~\ref{prop:erases_elimSpine_no_value} exists to refute; the
structural reason this chapter is longer than [S, Sect.~7.4] is that \emph{the pass is not a
congruence}. A hygiene point: \code{Lower.source\_app}, \code{Lower.source\_mkApps},
\code{Lower.target\_lambda} and \code{Lower.appReady} extend the \code{Lower} namespace from a
correctness module rather than from \code{Lower.lean}.
\end{remark}

\begin{definition}[Beta-readiness of a target head]
\label{def:BetaReady}
\lean{LeanToLambdaBox.BetaReady}
\leanok
\uses{def:WcbvEval}
$\mathsf{BetaReady}\ \Gamma\ fl\ X\ c$ says that the target reads $X$ as a function that beta-steps into
the body $c$: whenever a function position evaluates to $X$, an argument to $av$, and the instantiated
body $c[av]$ to $r$, the application evaluates to $r$. The Lean definition is \code{BetaReady}, stated
at arbitrary target flags.
\end{definition}

\begin{lemma}[The beta transport across the pass]
\label{lem:Lower-appReady}
\lean{LeanToLambdaBox.Lower.appReady, LeanToLambdaBox.Lower.target_lambda}
\leanok
\uses{def:BetaReady, def:Lower, def:ElimDecl}
At a specification environment with closed bodies and with guarded fix enabled, a lambda's image is
beta-ready at the image of the same body: if \Lower{} relates $\mathtt{lambda}\ nm\ b_0$ to $t$, there
is $c$ with \Lower{} relating $b_0$ to $c$ and $\mathsf{BetaReady}\ \Gamma\ fl\ t\ c$. The Lean
statement is \code{Lower.appReady}; \code{Lower.target\_lambda} records which target shapes a lambda can
take.
\begin{proof}
\leanok
\uses{def:LowerBlock, thm:Lower-constToFix, prop:LowerBlock-substList_fixSubst,
prop:LowerBlock-lambda_of_fixLambda, prop:Lower-closed}
Case analysis on the \Lower{} derivation at a lambda. Either the image is a lambda and the target
beta-reduces outright, or the lambda is a block member's own specification body, in which case the image
is the block's \code{fix} node and the target unfolds the fixpoint at principal argument index $0$
(guarded) to a lambda again --- beta-ready in turn, because a block member's emitted body is
lambda-headed. Every other arm is refuted on the node shape.
\end{proof}
\end{lemma}

\begin{remark}
The measured footprint of \code{Lower.appReady} is sorry-free, and free of choice as well; its row is
one of the forty in the committed fixture \code{test/ledger.expected}. This lemma is the only reason the
beta arm survives the pass's fix-closing; it is the analogue of the commutation of \code{tFix} unfolding
with beta on MetaRocq's emitted side, and it is where the convention of
Section~\ref{sec:correctness:closefix} is cashed in.
\end{remark}

\section{The environment accumulator and the emitted constructor arity}
\label{sec:correctness:env}

\begin{lemma}[The environment relation is antitone in the program]
\label{lem:ErasesEnv-subterm}
\lean{LeanToLambdaBox.ErasesEnv.subterm}
\leanok
\uses{def:ErasesEnv, def:SubTerm}
If \ErasesEnv{} holds of $t$ and $d$ is a subterm of $t$, then \ErasesEnv{} holds of $d$. The Lean
statement is \code{ErasesEnv.subterm}.
\begin{proof}
\leanok
\uses{lem:ReachableFrom-through_body, def:ErasesEnv}
Each of \ErasesEnv{}'s seven clauses requires something only of kernames \emph{reachable} from the term,
and a subterm reaches no more, so every clause transports. This is the most used lemma of the chapter:
each arm hands its induction hypothesis the environment relation at the subterm it recurses into. The
design point behind it is that \ErasesEnv{} is triggered by reachability of kernames, not by syntax,
which is what makes the closure kit below provable and hence what makes the induction go through.
\end{proof}
\end{lemma}

\begin{lemma}[The closure kit for the environment relation]
\label{lem:ErasesEnv-ofReach}
\lean{LeanToLambdaBox.ErasesEnv.ofReach, LeanToLambdaBox.ErasesEnv.box,
LeanToLambdaBox.ErasesEnv.substPair}
\leanok
\uses{def:ErasesEnv, def:ReachableFrom}
\ErasesEnv{} transports along any map that does not increase reachable kernames
(\code{ErasesEnv.ofReach}); $\Box$ reaches nothing, so it always satisfies the relation
(\code{ErasesEnv.box}); and the contractum $b[s]$ of a beta or zeta step names no kername its two parts
do not (\code{ErasesEnv.substPair}).
\begin{proof}
\leanok
\uses{lem:ErasesEnv-subterm, lem:ReachableFrom-through_body}
All three reduce to the reachability calculus: the first is immediate from the shape of the clauses, and
the other two exhibit the required reachability map ($\Box$ has none; a substitution's reachable set is
the union of the two parts').
\end{proof}
\end{lemma}

\begin{lemma}[The environment relation at a spine]
\label{lem:ErasesEnv-mkApps}
\lean{LeanToLambdaBox.ErasesEnv.mkApps}
\leanok
\uses{def:ErasesEnv, def:LBTerm-mkApps, def:ReachableFrom}
A spine is answered by whatever answers its head and its arguments: if \ErasesEnv{} holds of $f$ and of
every member of $l$, then it holds of $\mathsf{mkApps}\ f\ l$. The Lean statement is
\code{ErasesEnv.mkApps}.
\begin{proof}
\leanok
\uses{lem:ErasesEnv-ofReach, lem:ReachableFrom-through_body, def:ErasesEnv}
Re-establish each of the seven clauses at the spine, splitting the reachability witness by the inversion
for reachability through an application spine: a kername reached from the spine is reached from the head
or from one of the arguments, and the corresponding clause of that part answers it.
\end{proof}
\end{lemma}

\begin{remark}
Together with Lemma~\ref{lem:ErasesEnv-subterm} these are the accumulator lemmas that let the motive
carry \ErasesEnv{} at the \emph{value}; without them the induction hypothesis could not be applied at
any contractum. The capstone of Chapter~\ref{chap:capstone} spends \code{ErasesEnv.mkApps} again, at the
observable spine.
\end{remark}

\begin{lemma}[The emitted constructor arity of a reached block]
\label{lem:ErasesEnv-ctorArity}
\lean{LeanToLambdaBox.ErasesEnv.ctorArity}
\leanok
\uses{def:ErasesEnv, def:LowerEnv, def:IndInfo, def:constructorArity, def:ReachableFrom}
If \ErasesEnv{} holds of $t$, \LowerEnv{} prunes $\Gamma_{spec}$ to $\Gamma$, the source block data is
$\mathsf{IndInfo}\ env\ I\ iid\ np\ nfs$, the program reaches the block kername of $iid$, and
$k < |nfs|$, then the emitted environment reports constructor arity $np + nfs_k$ for constructor $k$ of
$iid$. The Lean statement is \code{ErasesEnv.ctorArity}.
\begin{proof}
\leanok
\uses{def:ErasesEnv, def:LowerEnv, def:ElimDecl}
Reachability of the block kername triggers \ErasesEnv{}'s \code{blocks} clause, which holds the block's
own emitted declaration together with its numbers; \LowerEnv{} says the pruning carries inductive blocks
over unchanged; unfolding the arity function on that declaration gives $np + nfs_k$, the number the
\lbox{} constructor-value rules read. The measured footprint is sorry-free.
\end{proof}
\end{lemma}

\begin{lemma}[A constructor spine within its arity evaluates componentwise]
\label{lem:wcbvEval_mkApps_construct}
\lean{LeanToLambdaBox.wcbvEval_mkApps_construct}
\leanok
\uses{def:WcbvEval, def:eraseFlags, def:constructorArity, def:LBTerm-mkApps}
If constructor $k$ of $iid$ has emitted arity $ar$ in $\Gamma$, $ts$ is no longer than $ar$, and each
$ts_i$ evaluates to $vs_i$, then $\mathsf{mkApps}\ (\mathtt{construct}\ iid\ k\ [\,])\ ts$ evaluates to
$\mathsf{mkApps}\ (\mathtt{construct}\ iid\ k\ [\,])\ vs$. The Lean statement is
\code{wcbvEval\_mkApps\_construct}, at \code{eraseFlags}.
\begin{proof}
\leanok
\uses{def:WcbvEval, def:eraseFlags}
Strong induction on the spine length, peeling the \emph{last} argument. The atomic constructor rule
starts the run at the bare declared constructor node and the applied constructor rule accumulates one
evaluated argument at a time, its saturation side condition supplied by the arity bound.
\end{proof}
\end{lemma}

\begin{remark}
This is the \emph{applied-form} constructor regime: \code{eraseFlags} sets
\code{with\_constructor\_as\_block} to false, the convention the Lean frontend emits, which differs from
the block form CertiRocq expects and agda2lambox produces. See Chapter~\ref{chap:lambdabox} for the flag
point and the rule shapes.
\end{remark}

\begin{lemma}[A constructor's index is in range]
\label{lem:CtorOf-lt_nfs}
\lean{LeanToLambdaBox.CtorOf.lt_nfs}
\leanok
\uses{def:UpstreamAsks, def:CtorOf, def:IndInfo}
Under \code{UpstreamAsks\ env}, if $c$ is constructor $k$ of $I$ and $I$'s block data is
$(iid, np, nfs)$, then $k < |nfs|$. The Lean statement is \code{CtorOf.lt\_nfs}.
\begin{proof}
\leanok
\uses{def:UpstreamAsks, lem:IndInfo-inj}
The witness inside \code{CtorOf} exhibits the member and the index; the declaration-level uniqueness the
upstream bundle asks for identifies the field-count list with the member's own constructor field counts,
whose length is that member's constructor count, and the index lookup \code{CtorOf} records would be
\code{none} above it.
\end{proof}
\end{lemma}

\section{An under-applied eliminator spine erases nothing that has a value}
\label{sec:correctness:elimspine}

This is the most intricate piece of engineering in the chapter, and its measured footprint is
sorry-free. The second \Lower{} reading of Lemma~\ref{lem:Lower-source_app} makes the beta rule's
\emph{function} a saturated eliminator spine one minor premise short. Refuting it means proving a
statement about the \emph{source} evaluator, reached from the \emph{target} side in three steps.

\begin{lemma}[No erasure image is a case node, a fix node, or an eliminator body]
\label{lem:erases_ne_case}
\lean{LeanToLambdaBox.erases_ne_case, LeanToLambdaBox.erases_ne_fix,
LeanToLambdaBox.erases_ne_elimBody}
\leanok
\uses{def:Erases, def:ElimBody}
If $e$ erases to $t$ then $t$ is not a \code{case} node, not a \code{fix} node, and does not satisfy
\code{ElimBody} for any $(iid, np, dp, nfs)$. The three Lean statements are \code{erases\_ne\_case},
\code{erases\_ne\_fix} and \code{erases\_ne\_elimBody}.
\begin{proof}
\leanok
\uses{def:Erases, def:ElimBody}
The source language has no match node and no fixpoint node --- block closure is the \emph{pass}'s job,
not the erasure's --- so no rule of \Erases{} emits either. \code{ElimBody} has two shapes, a
\code{case} under a lambda telescope or a \code{fix}; the telescope is peeled by the auxiliary inversion
for erasures of iterated lambdas, and both shapes then carry a node the erasure never emits.
\end{proof}
\end{lemma}

\begin{lemma}[A runtime eliminator key belongs to a casesOn constant]
\label{lem:ErasesEnv-runtimeKey_isCasesOn}
\lean{LeanToLambdaBox.ErasesEnv.runtimeKey_isCasesOn}
\leanok
\uses{def:ErasesEnv, def:ElimDecl, def:ConstOrigin, def:Kername, def:ReachableFrom}
If \ErasesEnv{} holds of $t$, the source name $c$ has \code{ConstOrigin}, the program reaches the
kername of $c$, and that kername is a runtime eliminator key of $\Gamma_{spec}$, then $c$ is a
\code{casesOn} name and has no compiler body. The Lean statement is
\code{ErasesEnv.runtimeKey\_isCasesOn}.
\begin{proof}
\leanok
\uses{lem:erases_ne_case, def:ErasesEnv}
Split on whether $c$ is tabled. If it is, \ErasesEnv{}'s \code{defns} clause says the entry at that key
is an erasure image, contradicting \code{erases\_ne\_elimBody} against the eliminator-body witness the
runtime key carries. If it is not, and $c$ were not a \code{casesOn} name, \ErasesEnv{}'s \code{axioms}
clause pins the entry to a body-less declaration, again contradicting the witness.
\end{proof}
\end{lemma}

\begin{remark}
The measured footprint is sorry-free. The doc-comment's own point is worth repeating: this is a
\emph{theorem} and not a clause of \ErasesEnv{} --- as a clause it would be false, since being a runtime
key tests the entry's body shape and never the key's name.
\end{remark}

\begin{lemma}[A source term whose erasure is a constant spine is itself one]
\label{lem:erases_constSpine_value}
\lean{LeanToLambdaBox.erases_constSpine_value, LeanToLambdaBox.erases_constSpine_head}
\leanok
\uses{def:Erases, def:SEval, def:ConstOrigin, def:Kername, def:LBTerm-mkApps}
If $e$ has a value and erases to $\mathsf{mkApps}\ (\mathtt{const}\ kn)\ args$, then there are $c$, $us$
and $as$ with the kername of $c$ equal to $kn$, with $c$ of \code{ConstOrigin}, with
$|as| \le |args|$, and with $\mkApps{}\ (\mathtt{const}\ c\ us)\ as$ evaluating to the same value. The
companion \code{erases\_constSpine\_head} says that in such an erasure the head erased by the constant
rule, so the kernames agree and the two spines have the same length.
\begin{proof}
\leanok
\uses{lem:Erases-sort_inv, def:Erases, def:SEval}
Induction on the evaluation. The literal arm unfolds a literal into constructor form; the beta arm
re-associates the head's spine with one more argument; the four spine arms are already at the shape;
every other arm's source erases to a node of another shape, which the head form and the inversion family
of \Erases{} refute. The produced spine is at most as long as the target's, which is all the consumer
reads.
\end{proof}
\end{lemma}

\begin{lemma}[A source eliminator spine below its arity heads no value]
\label{lem:SEval-no_elimSpine_value}
\lean{LeanToLambdaBox.SEval.no_elimSpine_value}
\leanok
\uses{def:SEval, def:UpstreamAsks, def:CasesOnShape, def:InformativeInd, def:ConstOrigin}
Let $c$ be a constant of \code{ConstOrigin} with no compiler body, such that every \emph{relevant}
\code{casesOn} shape at $c$ has $dp$ pre-arguments and $nm$ minor premises. Then no spine at $c$ with
fewer than $dp + 1 + nm$ arguments has an \SEval{} derivation. The Lean statement is
\code{SEval.no\_elimSpine\_value}.
\begin{proof}
\leanok
\uses{lem:consts_classified, def:SEval}
Strong induction on the spine length, then case analysis on every one of the eleven \SEval{} arms. The
delta arm is blocked by the absence of a compiler body; the constructor-value and type-former-value arms
by \code{ConstOrigin} excluding the other two classifications of a constant; the iota arm by the arity
agreement against the under-application bound; the beta arm by the induction hypothesis at the spine
with one argument removed; the remaining six arms by the subject's shape.
\end{proof}
\end{lemma}

\begin{proposition}[An under-applied eliminator spine erases nothing that has a value]
\label{prop:erases_elimSpine_no_value}
\lean{LeanToLambdaBox.erases_elimSpine_no_value}
\leanok
\uses{def:ErasesEnv, def:Erases, def:SEval, def:UpstreamAsks, def:ElimDecl, def:LBTerm-mkApps}
Under \code{UpstreamAsks\ env}, if \ErasesEnv{} holds of $\mathsf{mkApps}\ (\mathtt{const}\ kn)\ args$,
the key $kn$ carries \code{ElimDecl} data $(iid, np, dp, nfs)$, the spine is shorter than
$dp + 1 + |nfs|$, and $e$ erases to it, then $e$ has no \SEval{} value. The Lean statement is
\code{erases\_elimSpine\_no\_value}.
\begin{proof}
\leanok
\uses{lem:erases_constSpine_value, lem:ErasesEnv-runtimeKey_isCasesOn, lem:SEval-no_elimSpine_value,
lem:ErasesEnv-subterm, lem:ElimDecl-uniq, def:ErasesEnv}
The erasure's head is reached from the program, so Lemma~\ref{lem:ErasesEnv-runtimeKey_isCasesOn}
identifies the source head as a \code{casesOn} constant with no compiler body;
Lemma~\ref{lem:erases_constSpine_value} produces an \SEval{} derivation at that constant's spine;
\ErasesEnv{}'s \code{elims} clause turns any relevant \code{casesOn} shape at it into an
\code{ElimDecl}, which uniqueness of eliminator declarations equates with the entry's numbers, giving
exactly the arity agreement Lemma~\ref{lem:SEval-no_elimSpine_value} asks for; the under-application
bound closes it.
\end{proof}
\end{proposition}

\begin{remark}
The measured footprint is sorry-free. This is the one fact that makes the beta arm's second \Lower{}
reading refutable; without it T5 does not close.
\end{remark}

\section{The arms discharged in the aggregator}
\label{sec:correctness:arms}

\begin{lemma}[The type-former-value arm]
\label{lem:indVal_arm}
\lean{LeanToLambdaBox.indVal_arm}
\leanok
\uses{def:Simulates, def:IndInfo, def:UpstreamAsks}
An inductive-type-name spine --- a type former applied to arguments, a value in this weak evaluator ---
is simulated: its erasure can only be a boxed spine, and both sides box. The Lean statement is
\code{indVal\_arm}, taking the rule's own premises together with the induction hypothesis at each
argument.
\begin{proof}
\leanok
\uses{lem:erases_mkApps_inv, lem:erases_correct_boxLow, lem:ErasesEnv-subterm, lem:Erases-sort_inv,
lem:consts_classified, lem:CtorOf-not_indInfo, asm:lean4lean-trust}
The spine inversion offers the structural reading and the boxed-prefix reading. In the structural
reading the head would have to erase by the constructor rule or by the constant rule, and neither
applies to a type former; so only the boxed-prefix reading survives, and the spine form of
Lemma~\ref{lem:erases_correct_boxLow} closes it.
\end{proof}
\end{lemma}

\begin{remark}
The source rule exists because the delta and constructor-value arms evaluate their arguments and a
polymorphic call's type argument has exactly this shape --- worth saying, since a reader of [S] will not
expect a type-former-value rule.
\end{remark}

\begin{lemma}[The zeta (let) arm]
\label{lem:zeta_arm}
\lean{LeanToLambdaBox.zeta_arm}
\leanok
\uses{def:Simulates, def:ElimDecl, def:SEvalFlags}
A \code{let} whose bound value evaluates and whose instantiated body evaluates is simulated: either the
\code{let} boxes, or the emitted node is a \code{letIn} whose bound part evaluates by the induction
hypothesis and whose body's image is exactly the substitution instance. The Lean statement is
\code{zeta\_arm}; it asks for the zeta flag and for closed bodies in the specification environment.
\begin{proof}
\leanok
\uses{lem:erases_correct_boxLow, lem:ErasesEnv-ofReach, thm:SEval-defeq, lem:Erases-sort_inv,
prop:Lower-source_const, prop:Lower-subst_comm, thm:erases_subst, lem:Erases-defeqDFC_wt,
asm:lean4lean-trust}
After the box reading, the inversions of \Erases{} and of \Lower{} at a \code{let} give a \code{letIn}
node on both layers. The bound value's induction hypothesis evaluates its image; the body's image is the
substitution instance produced by the two substitution transports --- commutation of the pass with
substitution on the target side, the substitution lemma for \Erases{} on the erasure side --- after the
local context has been retyped, so that the body is translated against the \emph{evaluated} bound value.
The substitution closure of \ErasesEnv{} re-establishes the accumulator. This is the only arm that
manipulates a \code{let} entry of lean4lean's local context.
\end{proof}
\end{lemma}

\begin{proposition}[The beta arm]
\label{prop:beta_arm}
\lean{LeanToLambdaBox.beta_arm}
\leanok
\uses{def:Simulates, def:UpstreamAsks, def:ElimDecl, def:SEvalFlags}
An application whose function evaluates to a lambda, whose argument evaluates, and whose contractum
evaluates, is simulated. The Lean statement is \code{beta\_arm}; it asks for the beta flag, for
\code{UpstreamAsks\ env} and for closed bodies in the specification environment.
\begin{proof}
\leanok
\uses{lem:Lower-source_app, lem:Lower-appReady, prop:erases_elimSpine_no_value,
lem:erases_correct_boxLow, lem:ErasesEnv-ofReach, thm:SEval-defeq, lem:Erases-sort_inv,
prop:Lower-subst_comm, thm:erases_subst, prop:Lower-source_const, asm:lean4lean-trust}
After the box reading, the application node has two \Lower{} readings. In the congruence reading the
function's image is beta-ready by Lemma~\ref{lem:Lower-appReady} --- which covers both the lambda image
and the block \code{fix} image --- the argument's image evaluates by its induction hypothesis, and the
contractum's image is the substitution instance, \emph{unless} the function's own erasure is $\Box$, in
which case the whole application is irrelevant and both sides box. In the second reading the source is a
saturated eliminator spine one minor premise short, and the beta rule's own function premise is refuted
by Proposition~\ref{prop:erases_elimSpine_no_value}: such a spine heads no value.
\end{proof}
\end{proposition}

\begin{proposition}[The constructor-value arm]
\label{prop:ctorVal_arm}
\lean{LeanToLambdaBox.ctorVal_arm}
\leanok
\uses{def:Simulates, def:LowerEnv, def:CtorOf, def:IndInfo, def:UpstreamAsks}
A constructor spine within the constructor's arity, all of whose arguments evaluate, is simulated: its
image is the pointwise image of the arguments applied to the same \code{construct} node. The Lean
statement is \code{ctorVal\_arm}; the arity bound $|args| \le np + nfs_k$ is the rule's own premise.
\begin{proof}
\leanok
\uses{lem:erases_mkApps_inv, lem:ErasesEnv-ctorArity, lem:CtorOf-lt_nfs, lem:wcbvEval_mkApps_construct,
lem:Erases-mkApps_forall2, lem:ErasesEnv-mkApps, lem:erases_correct_boxLow, lem:Lower-source_app,
lem:Erases-sort_inv, lem:consts_classified, prop:Lower-source_const, lem:Lower-target_fix,
asm:lean4lean-trust}
The spine inversion's structural reading forces the head to erase to the \code{construct} node, the
constant rule being refuted by the rule's own \code{CtorOf}; the node's emitted arity is the one the
reached block declares, by Lemma~\ref{lem:ErasesEnv-ctorArity} with Lemma~\ref{lem:CtorOf-lt_nfs} for
the index bound; the pass is a congruence at a \code{construct}-headed spine; and the target value is
built argument by argument by Lemma~\ref{lem:wcbvEval_mkApps_construct}. The image of the value is then
the pointwise image of the arguments' values, assembled by Lemma~\ref{lem:Erases-mkApps_forall2}, by the
pass's own spine congruence and by Lemma~\ref{lem:ErasesEnv-mkApps}. The boxed-prefix reading folds as
usual.
\end{proof}
\end{proposition}

\begin{theorem}[T5 with the three hard arms as hypotheses]
\label{thm:erases_correct_of_steps}
\lean{LeanToLambdaBox.erases_correct_of_steps}
\leanok
\uses{def:ErasesCorrectStmt, def:StepIota}
Given inhabitants of \code{StepIota}, \code{StepProj} and \code{StepDelta}, the statement
$\code{ErasesCorrectStmt}\ env\ bo\ Us\ fl\ \Gamma_{spec}\ \Gamma$ holds. The Lean statement is
\code{erases\_correct\_of\_steps}.
\begin{proof}
\leanok
\uses{def:Simulates, lem:indVal_arm, lem:zeta_arm, prop:beta_arm, prop:ctorVal_arm,
lem:erases_correct_boxLow, lem:Erases-sort_inv, prop:Lower-source_const, def:LowerEnv,
asm:lean4lean-trust}
One induction on the \SEval{} derivation with motive \code{Simulates}, at the emitted environment and
with the local context specialised to the empty one immediately. The four trivial arms are discharged
inline --- a type erases to $\Box$ and evaluates to itself; a lambda is a value on both sides, whether
its image is a lambda or a block's \code{fix} node; a literal takes its one-step unfolding at the same
image. The four substantial non-step arms are the preceding nodes; the three remaining arms are the
hypotheses. Nothing beyond the eight binders of \code{ErasesCorrectStmt} is required: the forward
environment facts each arm needs are clauses of \ErasesEnv{} itself, and the kernel facts are fields of
\code{UpstreamAsks}.
\end{proof}
\end{theorem}

\begin{remark}
The aggregator closes with three non-vacuity witnesses --- \code{ctorVal\_target\_fires},
\code{indVal\_head\_boxes\_fires} and \code{beta\_reassociation\_fires} --- exhibiting respectively the
constructor spine the constructor-value arm builds, the head exclusion the type-former arm runs on, and
the re-associated pair the beta arm proceeds with at an over-applied eliminator spine. They are evidence
rather than results and get no node.
\end{remark}

\section{The three step arms: iota, delta, projection}
\label{sec:correctness:steps}

\begin{lemma}[Closed erasures, and the peeled alternative]
\label{lem:Erases-lbClosed}
\lean{LeanToLambdaBox.Erases.lbClosed, LeanToLambdaBox.LowerAlt.lower}
\leanok
\uses{def:Erases, def:LBClosed, def:Lower}
(a) If $e$ erases to $t$ in local context $\Delta$, then $t$ is \code{LBClosed} at $\Delta$'s
bound-variable count; in particular an erasure at the empty context is a closed \lbox{} term. (b) If the
alternative-peeling relation turns a minor premise $m$ into a \code{case} alternative
$(\mathit{names}, \mathit{body})$ for a constructor of $nf$ fields, then $|\mathit{names}| = nf$ and
\Lower{} relates $m$ to the lambda telescope over $\mathit{names}$ and $\mathit{body}$. The Lean
statements are \code{Erases.lbClosed} and \code{LowerAlt.lower}.
\begin{proof}
\leanok
\uses{def:Erases, def:Lower}
(a) Induction on the erasure: the variable rule copies only an index the local context answers, and the
two binder rules recurse under an extended context. (b) Induction on the peeling derivation, which
removes one binder at a time.
\end{proof}
\end{lemma}

\begin{remark}
Part (a) is needed because the iota bridge substitutes the constructor's fields simultaneously and in
reverse, which agrees with the beta chain only at closed field values; part (b) is what lets the iota
arm take its branch induction hypothesis at the \emph{un-peeled} alternative. \code{Erases.lbClosed} is
a general property of the erasure relation that happens to live in the iota arm's module, the iota
bridge being its first consumer.
\end{remark}

\begin{definition}[The two readings of a lowered constant-headed spine]
\label{def:LowerConstApp}
\lean{LeanToLambdaBox.LowerConstApp}
\leanok
\uses{def:Lower, def:ElimDecl, def:LBTerm-mkApps}
$\mathsf{LowerConstApp}\ \Gamma\ kn\ ts\ t$ holds when either (a) $kn$ is no runtime eliminator key and
$t$ is a congruence image of the spine --- head to head, argument to argument, lengths preserved --- or
(b) $kn$ carries \code{ElimDecl} data, the argument list splits as
$pre \mathbin{+\!\!+} disc :: minors \mathbin{+\!\!+} extra$ with $|pre| = dp$ and one minor premise per
constructor, each minor premise peels to an alternative, and
$t = \mathsf{mkApps}\ (\mathtt{case}\ (iid, np)\ disc'\ alts)\ extra'$ --- the over-application riding
\emph{outside} the emitted node. The Lean definition is \code{LowerConstApp}. Over-application costs no
premise, because the arguments past the node's arity are exactly the $extra$ of both the source iota
rule and the pass's eliminator rule.
\end{definition}

\begin{lemma}[Inverting the pass at a constant-headed spine]
\label{lem:Lower-source_constApp}
\lean{LeanToLambdaBox.Lower.source_constApp}
\leanok
\uses{def:LowerConstApp, def:Lower, def:LBTerm-mkApps}
If \Lower{} relates $\mathsf{mkApps}\ (\mathtt{const}\ kn)\ ts$ to $t$ then
$\mathsf{LowerConstApp}\ \Gamma\ kn\ ts\ t$. The Lean statement is \code{Lower.source\_constApp}, whose
induction is on a supplied spine length.
\begin{proof}
\leanok
\uses{lem:Lower-source_app, prop:Lower-source_const, lem:Lower-ne_fix_of_block}
Induction on the spine length. Argument by argument the derivation is either the application congruence
or the eliminator rule; an eliminator node at a shorter spine absorbs the remaining arguments into its
over-application. Both \code{fix}-producing arms are excluded by the block's own lambda-headedness: an
application is neither a constant nor a lambda.
\end{proof}
\end{lemma}

\begin{theorem}[The iota arm]
\label{thm:step_iota}
\lean{LeanToLambdaBox.step_iota}
\leanok
\uses{def:StepIota}
$\code{StepIota}\ env\ bo\ Us\ fl\ \Gamma_{spec}\ \Gamma$ holds: at a \code{casesOn} spine whose
discriminant reaches a constructor value, the emitted \code{case} node performs the same reduction,
over-application included. The Lean statement is \code{step\_iota}.
\begin{proof}
\leanok
\uses{lem:erases_mkApps_inv, lem:Lower-source_constApp, lem:Erases-lbClosed, lem:Erases-mkApps_forall2,
lem:erases_correct_boxLow, lem:ErasesEnv-ofReach, lem:ErasesEnv-subterm, thm:SEval-defeq,
lem:Lower-source_app, lem:ElimDecl-uniq, thm:not_erasable_of_informative, thm:elim_major,
thm:ctor_saturated, def:LowerEnv, prop:wcbvEval_mkApps_mkLambdas_substList,
lem:wcbvEval_mkApps_head_congr, lem:value_mkApps_construct_args, asm:lean4lean-trust}
Six steps. The spine inversion splits the erasure, the boxed-prefix case folding by the boxed-application
lemma. The head is the constant reading, and Lemma~\ref{lem:Lower-source_constApp} reads the pass back at
it; the congruence disjunct is refuted by the \code{ElimDecl} that \ErasesEnv{}'s \code{elims} clause
yields at the reached key, and uniqueness of eliminator declarations equates the node's data with that
one's. The discriminant's induction hypothesis gives a target constructor value, its boxed readings
refuted by irrelevance failing at an informative inductive, against the typing the major premise
supplies. The block's arity data comes through \LowerEnv{}, and the rule's own premise names the same
block. The branch's induction hypothesis is taken at the un-contracted application of the selected minor
premise to the constructor's fields (past the parameters) and the over-application. Finally the iota
bridge's beta-chain rewrite, under the head-congruence lemma for spines, turns that run into one iota
step at the emitted \code{case} node.
\end{proof}
\end{theorem}

\begin{remark}
The measured footprint of \code{step\_iota} contains \code{sorryAx}, and its row is in
\code{test/ledger.expected}. Non-vacuity witnesses live in the module's eliminator fixture:
\code{iota\_branch\_fires} and \code{iota\_overapplied\_fires} exhibit the emitted node firing with and
without over-application.
\end{remark}

\begin{lemma}[At a non-key constant head the pass is a congruence]
\label{lem:Lower-source_constSpine}
\lean{LeanToLambdaBox.Lower.source_constSpine}
\leanok
\uses{def:Lower, def:ElimDecl, def:LBTerm-mkApps}
If $kn$ is not a runtime eliminator key of $\Gamma$ and \Lower{} relates
$\mathsf{mkApps}\ (\mathtt{const}\ kn)\ ts$ to $t$, then $t$ is a congruence image: head to head,
argument to argument, lengths preserved. The Lean statement is \code{Lower.source\_constSpine}.
\begin{proof}
\leanok
\uses{lem:Lower-source_constApp, def:LowerConstApp}
The second disjunct of Definition~\ref{def:LowerConstApp} requires an \code{ElimDecl} at the key, which
the guard excludes; so the first disjunct holds.
\end{proof}
\end{lemma}

\begin{lemma}[What a tabled constant's image is]
\label{lem:Lower-const_body}
\lean{LeanToLambdaBox.Lower.const_body, LeanToLambdaBox.Lower.fixBody_of_block}
\leanok
\uses{def:Lower, def:ElimDecl, def:LowerBlock}
If \Lower{} relates $\mathtt{const}\ kn$ to $t$ and $\Gamma$ declares $kn$ with specification body $b$,
then $kn$ is no runtime eliminator key and either $t = \mathtt{const}\ kn$ or \Lower{} relates $b$ to
$t$. The supporting fact \code{Lower.fixBody\_of\_block} is that a block member's own body has the
block's \code{fix} node as a \Lower{} image, at the member the declaration pins.
\begin{proof}
\leanok
\uses{prop:LowerBlock-lambda_of_fixLambda, def:Lower}
Case analysis on the arms of \Lower{} that can produce a constant subject: the constant rule gives the
constant itself; the rule emitting a block from a constant gives the block's \code{fix} node, which
\code{Lower.fixBody\_of\_block} exhibits as an image of the constant's own body; the lambda-body and
eliminator arms are excluded on the node shape, and a runtime key would make the entry an eliminator
body.
\end{proof}
\end{lemma}

\begin{theorem}[The delta arm]
\label{thm:step_delta}
\lean{LeanToLambdaBox.step_delta}
\leanok
\uses{def:StepIota}
$\code{StepDelta}\ env\ bo\ Us\ fl\ \Gamma_{spec}\ \Gamma$ holds: unfolding a \emph{compiler} body at
the call site's universe levels is simulated, whether the emitted head is the constant's own body or the
block's \code{fix} node. The Lean statement is \code{step\_delta}.
\begin{proof}
\leanok
\uses{lem:erases_mkApps_inv, lem:Lower-source_constSpine, lem:Lower-const_body,
lem:Erases-mkApps_forall2, lem:ErasesEnv-ofReach, lem:ErasesEnv-mkApps, lem:erases_correct_boxLow,
lem:erases_ne_case, def:ErasesEnv, lem:ReachableFrom-through_body, def:LowerEnv,
prop:WcbvEval-mkApps_congr, lem:WcbvEval-head_value_of_mkApps, lem:consts_classified,
asm:lean4lean-trust}
The head erases to the tabled constant's kername; the constructor reading is refuted through
\ErasesEnv{}'s \code{tabled} clause and the classification of constants, and the boxed readings fold.
The spine lowers as a congruence, because a tabled constant is no runtime key: its specification body is
an erasure image, and no erasure image is an eliminator body. The induction hypothesis at the unfolded
application is then taken at the composite image whose \emph{head} is what the target head evaluates to
--- the emitted body when the pass relates the constant to itself, the block's \code{fix} node when it
relates it to the block --- and the spine congruence of \lbox{} evaluation moves the resulting run onto
the actual spine, replacing the head by one delta step or by nothing.
\end{proof}
\end{theorem}

\begin{remark}
The measured footprint contains \code{sorryAx}. The design point worth stating is that \emph{no fix
unfolding happens in this arm} --- a recursive constant's own run is the induction hypothesis's --- and
this is possible only because \code{Simulates} quantifies over \emph{all} composite images rather than a
computed one; that in turn is what Lemma~\ref{lem:Lower-appReady} pays for in the beta arm. The module's
witnesses \code{delta\_recursive\_fires} and \code{delta\_nonrecursive\_fires} exhibit both target
shapes.
\end{remark}

\begin{theorem}[The projection arm]
\label{thm:step_proj}
\lean{LeanToLambdaBox.step_proj}
\leanok
\uses{def:StepIota}
$\code{StepProj}\ env\ bo\ Us\ fl\ \Gamma_{spec}\ \Gamma$ holds: a structure projection whose
discriminant reaches a constructor spine is simulated by the emitted \code{proj} node, which carries the
same triple of inductive identifier, parameter count and field index. The Lean statement is
\code{step\_proj}.
\begin{proof}
\leanok
\uses{lem:erases_mkApps_inv, lem:erasable_mkApps, lem:erases_correct_boxLow, lem:CtorOf-lt_nfs,
thm:SEval-defeq, lem:Lower-source_app, lem:Erases-sort_inv, prop:Lower-source_const, def:ErasesEnv,
def:LowerEnv, thm:not_erasable_of_informative, asm:lean4lean-trust}
The shortest of the three, because a projection lowers to a projection and there is no spine to invert.
After the box reading, the erasure is a \code{proj} node whose block data is the source block's; the
discriminant's induction hypothesis gives the target constructor value, its boxed readings refuted by
irrelevance failing at an informative inductive, against the typing lean4lean's projection judgement
carries; the node's non-propositional flag comes from \ErasesEnv{}'s \code{blocks} clause through
\LowerEnv{}; the field's induction hypothesis is taken at the selected argument, and the target's
projection rule assembles. The prefix the node drops is the parameter count the source rule names, and
the constructor index is forced to $0$ at a single-constructor block.
\end{proof}
\end{theorem}

\begin{remark}
The measured footprint contains \code{sorryAx}. An honesty item the module records itself: the arm is
non-vacuous today only on hand-built witnesses, because its typing premise at a projection is
inhabitable only through lean4lean's projection judgement, whose kernel adequacy is \code{sorry} at the
pin. See Chapter~\ref{chap:trust}.
\end{remark}

\section{Closing the simulation}
\label{sec:correctness:close}

\begin{theorem}[T5: erasure correctness]
\label{thm:erases_correct}
\lean{LeanToLambdaBox.erases_correct}
\leanok
\uses{def:ErasesCorrectStmt}
For a well-formed kernel environment $env$, a closed Lean term $e$ that translates to $ve$ and evaluates
to $v$ under the source evaluator, and any \lbox{} program $t$ obtained by erasing $e$ to $t_0$ and
lowering $t_0$ at a specification environment $\Gamma_{spec}$ that satisfies \ErasesEnv{} for $t_0$ and
whose pruning to $\Gamma$ satisfies \LowerEnv{} --- and granted the four upstream kernel facts
\code{UpstreamAsks\ env} --- the emitted program evaluates $t$, under \lbox{}'s weak call-by-value
semantics at \code{eraseFlags} (prop-case off, guarded fix, applied-form constructors), to a term $v'$
that is itself an erasure-then-lowering image of the source value $v$. The Lean statement is
\code{erases\_correct}, which \emph{is} Definition~\ref{def:ErasesCorrectStmt}: the source flag record
stays a parameter, so the theorem is unconditional in the source flags.
\begin{proof}
\leanok
\uses{thm:erases_correct_of_steps, thm:step_iota, thm:step_proj, thm:step_delta, asm:lean4lean-trust}
Theorem~\ref{thm:erases_correct_of_steps} applied to the three arm theorems; one line, the content being
in the nodes above.
\end{proof}
\end{theorem}

\begin{remark}
The measured footprint of \code{erases\_correct} contains \code{sorryAx}, inherited from the pinned
lean4lean and not introduced here; its row heads the committed fixture \code{test/ledger.expected}. The
theorem is applied exactly once downstream, at the observable spine inside the capstone
(Chapter~\ref{chap:capstone}). It is conditional --- on \code{UpstreamAsks\ env} and on the two
environment premises --- and must never be quoted without them.
\end{remark}

\begin{theorem}[T5, folded through the composite relation]
\label{thm:erases_correct_lb}
\lean{LeanToLambdaBox.erases_correct_lb}
\leanok
\uses{def:ErasesCorrectLBStmt, def:ErasesLB}
$\code{ErasesCorrectLBStmt}\ env\ bo\ Us\ fl\ \Gamma_{spec}\ \Gamma$ holds: the same simulation read
through \code{ErasesLB}, with the environment premise taken at whichever middle term the composite
exhibits. The Lean statement is \code{erases\_correct\_lb}.
\begin{proof}
\leanok
\uses{thm:erases_correct, def:ErasesLB, asm:lean4lean-trust}
Destructure the \code{ErasesLB} witness into its erasure and lowering halves, apply
Theorem~\ref{thm:erases_correct} to them and to the environment premise instantiated at the exhibited
middle term, and repackage the conclusion as an \code{ErasesLB} witness.
\end{proof}
\end{theorem}

\section{The fix-closing convention}
\label{sec:correctness:closefix}

A short module pins the de Bruijn convention shared by the shipping eraser's block-closing loop and
\lbox{}'s \code{fix} node. It mentions no source-side notion at all.

\begin{definition}[The n-way abstraction of a block's fresh variables]
\label{def:closeFix}
\lean{LeanToLambdaBox.closeFix, LeanToLambdaBox.closeFixFold}
\leanok
\uses{def:toBvar, def:LBTerm}
\code{closeFixFold} folds the single-variable abstraction \code{toBvar} over a list of (free variable,
level) instructions, head first; $\code{closeFix}\ \mathit{ids}\ \mathit{base}\ t$ runs it on
$\mathit{ids}$ reversed and indexed from $\mathit{base}$, so that for a block
$\mathit{ids} = [x_0,\dots,x_{m-1}]$ the \emph{last} sibling $x_{m-1}$ becomes
$\mathtt{bvar}\ \mathit{base}$ and the \emph{first} $x_0$ becomes
$\mathtt{bvar}\ (\mathit{base}+m-1)$ --- exactly the convention of the \lbox{} \code{fix} node, whose
bodies live under $m$ binders and refer to sibling $k$ by $\mathtt{bvar}\ (m-1-k)$.
\end{definition}

\begin{lemma}[The metatheory's recursion is the shipping loop]
\label{lem:closeFixFold_eq_foldl}
\lean{LeanToLambdaBox.closeFixFold_eq_foldl}
\leanok
\uses{def:closeFix, def:toBvar}
$\code{closeFixFold}\ \mathit{pairs}\ t$ equals the left fold of \code{toBvar} over $\mathit{pairs}$
starting from $t$. The Lean statement is \code{closeFixFold\_eq\_foldl}.
\begin{proof}
\leanok
\uses{def:closeFix}
Induction on the instruction list, generalising the accumulator.
\end{proof}
\end{lemma}

\begin{remark}
This is the bridge from the shipping code's imperative loop --- which elaborates to exactly this left
fold --- to the metatheory's structural recursion. What is \emph{not} proved here is the \code{fixvars}
lookup layer, \ie{} that the variable recorded for member $k$ is the one the closing instruction names;
that is a fact about the run, carried by the run model and the block-lowering relation of
Chapters~\ref{chap:run} and~\ref{chap:lowering}.
\end{remark}

\begin{lemma}[Per-sibling abstraction, occurrence reflection, and a worked two-member block]
\label{lem:closeFixFold_fvar_head}
\lean{LeanToLambdaBox.closeFixFold_fvar_head, LeanToLambdaBox.hasFVar_closeFix_of,
LeanToLambdaBox.closeFix_not_hasFVar, LeanToLambdaBox.closeFix_2block_last,
LeanToLambdaBox.closeFix_2block_first}
\leanok
\uses{def:closeFix, def:LBTerm}
If $x$ is closed at level $\mathit{lvl}$ by the first instruction then
$\code{closeFixFold}\ ((x,\mathit{lvl}) :: \mathit{rest})\ (\mathtt{fvar}\ x) =
\mathtt{bvar}\ \mathit{lvl}$, regardless of the remaining instructions --- once $x$ is a \code{bvar},
every subsequent abstraction is a no-op. Closing removes exactly the block identifiers and introduces
none: a variable occurring after closing occurred before and is none of the abstracted ones, and a
variable that either did not occur or is abstracted does not occur after. Worked out on a two-member
block, the last sibling becomes $\mathtt{bvar}\ 0$ and the first $\mathtt{bvar}\ 1$, so a body
$x_0\,x_1$ closes to $\mathtt{bvar}\ 1\ \mathtt{bvar}\ 0$ --- what the \code{fix} node expects.
\begin{proof}
\leanok
\uses{lem:closeFixFold_eq_foldl, def:closeFix}
All by induction on the instruction list, using that \code{toBvar} is the identity on a term with no
occurrence of its variable and introduces no free variable. The two-member statements are then instances
computed by unfolding.
\end{proof}
\end{lemma}

\begin{remark}
The convention is made a theorem plus a worked example rather than a comment because the whole
target-side fix machinery depends on it. Most of the module is nevertheless orphaned: only
\code{closeFix}, \code{closeFixFold} and the cons equation have consumers outside it, and the occurrence
results are independently re-proved in the block-lowering module.
\end{remark}

\section{The first-order fragment and the observation}
\label{sec:correctness:firstorder}

T5 says that \emph{whatever} the relation produces is reached; it does not say the program computes the
right answer, because the relation is non-deterministic. Following [S, Sect.~7.3], the non-determinism
collapses on values of first-order inductive types --- ``values of first-order types are just nested
constructor applications where neither proofs nor types can occur'' --- and that collapse is the
\emph{observation} the capstone reports.

\begin{definition}[A first-order block]
\label{def:FirstOrderDecl}
\lean{LeanToLambdaBox.FirstOrderDecl, LeanToLambdaBox.FOType}
\leanok
\uses{def:lean4lean-grounding}
A field type is \emph{first-order with respect to a closure} $fo$ and a block's own names when it is a
\textbf{bare} type former $\mathtt{const}\ I$ with $I$ among the block's own names or accepted by $fo$
(\code{FOType}); there is deliberately no bound-variable clause, since a constructor type is closed and
a bound variable at binder $i$ is a parameter or an earlier field, never a type former. A block
declaration is first-order when it satisfies four clauses: \code{mono} (no universe parameters),
\code{informative} (every member's declared result sort is a successor), \code{noIndices} (every
member's arity equals the parameter count) and \code{fields} (every binder of every constructor,
parameters included, is a first-order field type). The Lean structure is \code{FirstOrderDecl}.
\end{definition}

\begin{remark}
The clauses should be booked honestly. \code{informative} is the result-sort half of [L, Def.~6] in the
successor shape and \code{fields} is [L, Def.~14] restricted to unapplied field types, while
\code{mono} and \code{noIndices} are declared scope restrictions of this development --- the module says
so --- booked neither to Letouzey nor to MetaRocq's \code{firstorder\_ind}. The module header also
records that MetaRocq's shipped \code{firstorder\_ind} evaluates to false on \code{nat}, so this
development cites rather than transcribes it.
\end{remark}

\begin{definition}[First-order inductive types]
\label{def:FirstOrderInd}
\lean{LeanToLambdaBox.FirstOrderInd, LeanToLambdaBox.FOClosed, LeanToLambdaBox.HasInduct}
\leanok
\uses{def:FirstOrderDecl, def:lean4lean-grounding}
A set of names $fo$ is \emph{closed} when every name it accepts is declared by a block of $env$'s own
declaration list, all of whose members are first-order with respect to $fo$ itself (\code{FOClosed},
with \code{HasInduct} for ``declared by a block of $env$''). Then $I$ is a \emph{first-order inductive
type of $env$} when some closed $fo$ contains $I$. The Lean definition is \code{FirstOrderInd}; the
closure is existentially quantified \emph{inside} it rather than left a parameter, which the module
header documents as a soundness fix: with a free $fo$ the predicate is satisfied by the constantly-true
set, under which a type former with an arbitrary field would count as first-order.
\end{definition}

\begin{lemma}[A first-order type former is declared and informative]
\label{lem:FirstOrderInd-indDeclOf}
\lean{LeanToLambdaBox.FirstOrderInd.indDeclOf, LeanToLambdaBox.FirstOrderInd.informativeInd}
\leanok
\uses{def:FirstOrderInd, def:IndDeclOf, def:InformativeInd}
If $I$ is first-order in $env$ then $I$ belongs to a declared block and $I$ is informative. The Lean
statements are \code{FirstOrderInd.indDeclOf} and \code{FirstOrderInd.informativeInd}.
\begin{proof}
\leanok
\uses{lem:InformativeInd-mono, lem:IndInfo-constant_isArity, def:FirstOrderDecl}
Unfold the closure at $I$ to get the declaring block, then read declaredness off it; informativity
follows from the \code{informative} clause's successor shape, whose level never evaluates to zero, which
is the criterion informativity uses. This is the hinge between the \emph{syntactic} first-order
condition and the \emph{semantic} relevance predicate the target-facing relations read.
\end{proof}
\end{lemma}

\begin{lemma}[A first-order spine is neither a sort nor a Pi]
\label{lem:FirstOrderInd-notSortNotPi}
\lean{LeanToLambdaBox.FirstOrderInd.notSortNotPi}
\leanok
\uses{def:FirstOrderInd, def:UpstreamAsks, def:lean4lean-grounding}
Under \code{UpstreamAsks\ env}, in a well-formed context, if $I$ is first-order and
$\mathsf{mkApps}\ (\mathtt{const}\ I\ ius)\ iargs$ is typeable, then that spine is definitionally equal
to no sort and to no Pi type. The Lean statement is \code{FirstOrderInd.notSortNotPi}.
\begin{proof}
\leanok
\uses{lem:FirstOrderInd-indDeclOf, def:UpstreamAsks}
The declaring block is exhibited by the closure, so the upstream ask about the arities of declared
constants applies at it.
\end{proof}
\end{lemma}

\begin{remark}
This is what excludes lambda values and type values from the first-order induction; it rests on one of
the four fields of \code{UpstreamAsks}, so it is a \emph{hypothesis} of the development and not a
theorem about lean4lean.
\end{remark}

\begin{definition}[The Boolean first-order checker]
\label{def:firstOrderIndB}
\lean{LeanToLambdaBox.firstOrderIndB, LeanToLambdaBox.foMemberB, LeanToLambdaBox.foFieldB,
LeanToLambdaBox.piDomains}
\leanok
\uses{def:Witness-SourceTable}
A fuelled Boolean decision of the first-order condition on a \emph{reified} source table: $I$ is tabled,
belongs to the block it names, and every member of that block is monomorphic, index-free, lands in a
syntactically successor sort, and has only bare-type-former constructor binders, recursively at one less
unit of fuel. The Lean definitions are \code{firstOrderIndB}, with \code{foMemberB}, \code{foFieldB} and
\code{piDomains} for the member, field and binder-list layers.
\end{definition}

\begin{remark}
The module decides seven instances by reduction on a reified table (\code{Nat}, \code{Bool}, a benchmark
tree, a function-field fixture, \code{True}, \code{List}, \code{Prod}): the first three are true; the
function-field fixture is false, because \code{FOType} admits only bare formers; \code{True} is false,
because its result sort is not a successor; and \code{List} and \code{Prod} are false --- the module's
own comment says that \code{List Nat} and \code{Nat} $\times$ \code{Nat} are first-order data, but the
predicate is chosen at the \emph{declaration} level, so a universe-polymorphic former is out of the
domain. These seven facts are anonymous \code{example}s and therefore carry no citable Lean name; they
are reported here in prose only. The relevance test used here is the stricter syntactic-successor one
matching \code{FirstOrderDecl.informative}, not the never-zero test matching semantic informativity.
\end{remark}

\begin{lemma}[One step of the checker, read off the table]
\label{lem:firstOrderIndB_step}
\lean{LeanToLambdaBox.firstOrderIndB_step}
\leanok
\uses{def:firstOrderIndB, def:Witness-SourceTable}
If the checker returns true at fuel $\mathit{fuel}+1$ for $I$, then $I$ is tabled with a block value,
$I$ is among that block's members, and every member $J$ of the block is tabled with the same block, has
no level parameters, no indices, a successor result sort, and every constructor binder of the form
$\mathtt{const}\ K\ us$ with $K$ in the block or accepted by the checker at one less fuel. The Lean
statement is \code{firstOrderIndB\_step}.
\begin{proof}
\leanok
\uses{def:firstOrderIndB}
Unfold one layer of the fuelled recursion and convert each Boolean conjunct into its propositional
reading, using the checker's monotonicity in the fuel to widen the recursive call.
\end{proof}
\end{lemma}

\begin{remark}
This is the honesty item of the section: \code{firstOrderIndB\_step} is the \emph{table-side} half of
the checker's soundness only. What is missing --- and what the module header states explicitly --- is
that the tabled block is a block of the \emph{model} environment; that is one of the outstanding
upstream asks (inverting lean4lean's environment translation at an inductive name). The reduction
decisions of Definition~\ref{def:firstOrderIndB} therefore do \emph{not} establish
$\code{FirstOrderInd}\ env\ \code{Nat}$, and must not be read as doing so.
\end{remark}

\begin{definition}[The shapes a source evaluation returns]
\label{def:SValue}
\lean{LeanToLambdaBox.SValue}
\leanok
\uses{def:CtorOf, def:IndInfo}
An inductive predicate with five arms: a lambda, a sort, a Pi type, a constructor-headed spine of
values, and a type-former-headed spine of values. The Lean definition is \code{SValue}.
\end{definition}

\begin{remark}
\code{SValue} is strictly wider than the set of values \SEval{} returns: its constructor arm omits the
arity bound $|args| \le np + nfs_k$ that the source constructor-value rule carries --- the bound
[S, Fig.~12] states as $\mathit{nargs} \le \mathit{cstr\_arity}$, and without which the arm would call a
value a spine the target is stuck on. Lemma~\ref{lem:SEval-svalue} is still provable, but the
doc-comment's ``what an \code{SEval} derivation returns'' overstates the correspondence: it is a
superset.
\end{remark}

\begin{lemma}[A source evaluation returns a value]
\label{lem:SEval-svalue}
\lean{LeanToLambdaBox.SEval.svalue}
\leanok
\uses{def:SValue, def:SEval}
If \SEval{} relates $e$ to $v$ then $\mathsf{SValue}\ env\ v$. The Lean statement is
\code{SEval.svalue}.
\begin{proof}
\leanok
\uses{def:SValue, def:SEval}
One arm per rule: the three value arms conclude directly, and every reduction arm reads its result off
its continuation premise.
\end{proof}
\end{lemma}

\begin{lemma}[Every member of a first-order block is first-order]
\label{lem:firstOrderInd_of_own}
\lean{LeanToLambdaBox.firstOrderInd_of_own}
\leanok
\uses{def:FirstOrderInd, def:FirstOrderDecl}
If $fo$ is closed, $\mathit{decl}$ is a block of $env$ that is first-order with respect to $fo$, and $J$
is one of its members' names, then $J$ is first-order in $env$. The Lean statement is
\code{firstOrderInd\_of\_own}.
\begin{proof}
\leanok
\uses{def:FirstOrderDecl, def:FirstOrderInd}
Widen the closure by the block's own names. The closure occurs only \emph{positively} in
\code{FirstOrderDecl}, so widening preserves the post-fixed-point property, and the widened set contains
$J$.
\end{proof}
\end{lemma}

\begin{proposition}[Fields of a first-order constructor value are first-order]
\label{prop:fOFields_of_asks}
\lean{LeanToLambdaBox.fOFields_of_asks}
\leanok
\uses{def:FirstOrderInd, def:UpstreamAsks, def:CtorOf, def:lean4lean-grounding}
In a well-formed environment, under \code{UpstreamAsks\ env}: if $I$ is first-order, $c$ is a
constructor of some $I'$, the spine $\mkApps{}\ (\mathtt{const}\ c\ us)\ \mathit{cargs}$ translates to a
model term typed at a spine of $I$, and argument $i$ translates to $a$, then there is a first-order $J$
with $a$ typed at a spine of $J$. The Lean statement is \code{fOFields\_of\_asks}.
\begin{proof}
\leanok
\uses{lem:firstOrderInd_of_own, def:FirstOrderDecl, lem:trExprS_spine_head, lem:IndInfo-arity,
thm:ctor_saturated, lem:peel_piSpine, thm:elim_major, lem:IndInfo-inj, def:UpstreamAsks,
asm:lean4lean-trust}
Two fields of \code{UpstreamAsks} carry it: the constructor's own type former is the value's, against
the result spine that peeling the constructor type reaches; and the block the closure exhibits is the
block the constructor's classification exhibits, by declaration-level uniqueness, which is what lets the
\code{fields} clause name the field's former. The spine's typing then peels without a further ask,
because a \emph{translated} spine carries its argument typings in the translation's own application arm.
\end{proof}
\end{proposition}

\begin{remark}
The measured footprint contains \code{sorryAx}. This proposition is the recursion premise of
Theorem~\ref{thm:firstorder_erases_core}: it is what makes the induction on \code{SValue} both
well-founded and type-correct at each field.
\end{remark}

\begin{definition}[The lambda-box image of a first-order value]
\label{def:FOSpine}
\lean{LeanToLambdaBox.FOSpine}
\leanok
\uses{def:LBTerm}
An applied-form constructor tree: a bare \code{construct} node is one, and an application of two
constructor trees is one. The Lean definition is \code{FOSpine}, and those two rules are the whole
predicate, because a \code{construct} node carries \emph{no} arguments on the emitted output: the tree
is built by application. A non-vacuity \code{example} in the module exhibits the \lbox{} Peano numeral
$1$; being anonymous, it is cited here in prose only.
\end{definition}

\begin{lemma}[Constructor trees: box-freedom, spines, and transport across the pass]
\label{lem:FOSpine-lower}
\lean{LeanToLambdaBox.FOSpine.lower, LeanToLambdaBox.FOSpine.noBox, LeanToLambdaBox.FOSpine.mkApps,
LeanToLambdaBox.FOSpine.spineHead}
\leanok
\uses{def:FOSpine, def:Lower, def:NoBox, def:LBTerm-mkApps}
A constructor tree holds no box; a bare \code{construct} node applied to constructor trees is one; its
spine head is its \code{construct} node; and a constructor tree has a constructor tree for a lowered
image. The four Lean statements are \code{FOSpine.noBox}, \code{FOSpine.mkApps},
\code{FOSpine.spineHead} and \code{FOSpine.lower}.
\begin{proof}
\leanok
\uses{lem:Lower-source_app, prop:Lower-source_const, lem:Lower-ne_fix_of_block, def:FOSpine}
The first three are inductions on the tree. For the transport: the constructor arm goes through the
pass's inversion at a \code{construct} node, which excludes both \code{fix}-producing arms itself (a
\code{construct} node is neither a constant nor a lambda) and returns a zero-length argument list; the
application arm goes through Lemma~\ref{lem:Lower-source_app}, whose eliminator disjunct is refuted on
the head --- a constant there, a constructor here.
\end{proof}
\end{lemma}

\begin{lemma}[Box-freedom of the lowered first-order value]
\label{lem:noBox_lower_of_foSpine}
\lean{LeanToLambdaBox.noBox_lower_of_foSpine}
\leanok
\uses{def:FOSpine, def:Lower, def:NoBox}
If $t_{v_0}$ is a constructor tree and \Lower{} relates it to $t_v$, then $t_v$ satisfies \code{NoBox}.
The Lean statement is \code{noBox\_lower\_of\_foSpine}.
\begin{proof}
\leanok
\uses{lem:FOSpine-lower}
Transport the shape across the pass by Lemma~\ref{lem:FOSpine-lower}, then read box-freedom off it.
\end{proof}
\end{lemma}

\begin{remark}
Why it is stated at the \emph{shape} and not at \code{NoBox}: box-freedom does not transport along
\Lower{}. The arm that emits a block from a constant relates a box-free constant to a block whose
definitions carry their own boxes, and the repository records the counterexample. This lemma supplies
the box-freedom conjunct of the capstone's conclusion (Chapter~\ref{chap:capstone}).
\end{remark}

\begin{theorem}[The erasure of a first-order value is unique and a constructor tree]
\label{thm:firstorder_erases_core}
\lean{LeanToLambdaBox.firstorder_erases_core}
\leanok
\uses{def:SValue, def:FirstOrderInd, def:FOSpine, def:Erases, def:UpstreamAsks,
def:lean4lean-grounding}
In a well-formed environment, under \code{UpstreamAsks\ env}: for every source value $v$, every
first-order $I$, every translation of $v$ typed at a spine of $I$, and every erasure $t$ of $v$, the
term $t$ is a constructor tree, and every erasure of $v$ equals $t$. The Lean statement is
\code{firstorder\_erases\_core}.
\begin{proof}
\leanok
\uses{lem:FirstOrderInd-notSortNotPi, lem:FirstOrderInd-indDeclOf, prop:fOFields_of_asks,
lem:erases_mkApps_inv, lem:erasable_mkApps, lem:FOSpine-lower, lem:trExprS_spine_head,
thm:not_erasable_of_informative, lem:Erases-sort_inv, lem:Erases-sort_erasable,
lem:Erases-indInfo_erasable, lem:consts_classified, asm:lean4lean-trust}
One induction over the value's shape. A lambda is excluded because its type is a Pi and a first-order
spine is not (Lemma~\ref{lem:FirstOrderInd-notSortNotPi}); a sort, a Pi type and a type-former spine are
excluded because they are \Erasable{} while a first-order value is not, informativity being read off
Lemma~\ref{lem:FirstOrderInd-indDeclOf}. What remains is a constructor spine, which erases by the
congruence alone: its boxed readings fall to the same irrelevance fact, its head to the classification
of constants, and its arguments to the induction hypothesis at the field typings
Proposition~\ref{prop:fOFields_of_asks} supplies. Uniqueness of the whole spine then follows from
pointwise uniqueness.
\end{proof}
\end{theorem}

\begin{remark}
This is the single result that makes the capstone's observable meaningful: it delivers both uniqueness
of the answer and, through \code{FOSpine}, its box-freedom. Its measured footprint contains
\code{sorryAx}.
\end{remark}

\begin{theorem}[First-order erasure is deterministic]
\label{thm:firstorder_erases_deterministic}
\lean{LeanToLambdaBox.firstorder_erases_deterministic}
\leanok
\uses{def:FirstOrderInd, def:Erases, def:SEval, def:UpstreamAsks, def:lean4lean-grounding}
At a source \emph{value} of a first-order inductive type --- a term that evaluates to itself --- the
erasure relation has exactly one image: if $v$ erases to $t_1$ and to $t_2$ then $t_1 = t_2$. The Lean
statement is \code{firstorder\_erases\_deterministic}, whose self-evaluation premise is an \SEval{}
derivation from $v$ to $v$.
\begin{proof}
\leanok
\uses{thm:firstorder_erases_core, lem:SEval-svalue, asm:lean4lean-trust}
Lemma~\ref{lem:SEval-svalue} turns the self-evaluation into an \code{SValue} witness, and
Theorem~\ref{thm:firstorder_erases_core} then gives uniqueness.
\end{proof}
\end{theorem}

\begin{remark}
The measured footprint contains \code{sorryAx}. This is the point at which this development's relational
erasure recovers the functional character of MetaRocq's \code{erase} --- the analogue of [S]'s
\code{firstorder\_erases\_deterministic} (p.~8:63), whose conclusion is stated against the erasure
\emph{function} MetaRocq has and this development does not. It is not instantiated at \code{Nat}
anywhere in the tree; see Lemma~\ref{lem:FOModel-firstOrderInd_E}.
\end{remark}

\begin{theorem}[The erasure of a first-order value holds no box]
\label{thm:firstorder_no_box}
\lean{LeanToLambdaBox.firstorder_no_box}
\leanok
\uses{def:FirstOrderInd, def:Erases, def:NoBox, def:SEval, def:UpstreamAsks}
Same binders as the preceding theorem; if $v$ erases to $t$ then $t$ satisfies \code{NoBox}. The Lean
statement is \code{firstorder\_no\_box}.
\begin{proof}
\leanok
\uses{thm:firstorder_erases_core, lem:SEval-svalue, lem:FOSpine-lower, asm:lean4lean-trust}
Theorem~\ref{thm:firstorder_erases_core} returns the shape, and box-freedom is read off it.
\end{proof}
\end{theorem}

\begin{remark}
The measured footprint contains \code{sorryAx}. The caveat bears repeating: this is box-freedom of the
\emph{erasure}, not of the lowered image; the lowered statement is
Lemma~\ref{lem:noBox_lower_of_foSpine}, and that is the one the capstone uses.
\end{remark}

\begin{lemma}[FirstOrderInd is inhabited]
\label{lem:FOModel-firstOrderInd_E}
\lean{LeanToLambdaBox.FOModel.firstOrderInd_E, LeanToLambdaBox.FOModel.decl_wf,
LeanToLambdaBox.FOModel.wf'}
\leanok
\uses{def:FirstOrderInd, def:lean4lean-grounding}
There is a well-formed model environment in which $\code{FirstOrderInd}\ env\ E$ holds, for a hand-built
inductive $E$: empty (no constructors), \code{Type}-valued, parameter-free, index-free, with a single
\code{Prop}-eliminating recursor. The Lean statements are \code{FOModel.decl\_wf}, which discharges
every clause of lean4lean's well-formedness for the declaration, \code{FOModel.wf'}, which exhibits the
one-declaration well-formed list, and \code{FOModel.firstOrderInd\_E}.
\begin{proof}
\leanok
\uses{def:FirstOrderInd, def:FirstOrderDecl}
\code{decl\_wf} is discharged clause by clause, most of them vacuously since there are no constructors;
\code{firstOrderInd\_E} then exhibits the closure that accepts exactly $E$, whose \code{FirstOrderDecl}
obligations are vacuous for the same reason.
\end{proof}
\end{lemma}

\begin{remark}
This is the vacuity caveat of the whole section. It is the \emph{only} proved instance of
\code{FirstOrderInd} in the tree, and it is at an inductive with no constructors, so the substantive
branch of Theorem~\ref{thm:firstorder_erases_core} --- the constructor arm, which is its entire content
--- is vacuous at the only available instance. The module header says as much, and the trust document
records that all eight green rungs \emph{assume} $\code{FirstOrderInd}\ env\ \code{Nat}$ as a hypothesis
discharged nowhere; see Chapters~\ref{chap:capstone} and~\ref{chap:trust}.
\end{remark}

\section{Correspondence with the paper, and where this is weaker}
\label{sec:correctness:paper}

\begin{center}
\begin{tabular}{ll}
\textbf{[S]} & \textbf{here} \\
\code{simple\_erases\_correct} (p.~8:62, stated not proved) & --- \\
\code{erases\_correct} with \code{erases\_deps} (Sect.~7.4) & \code{erases\_correct} \\
$\Sigma ; [\,] \vdash t \leadsto t'$ & $\Erases{}\ env\ Us\ [\,]\ e\ t_0$ \\
$\Sigma \leadsto \Sigma'$ / \code{erases\_deps} & $\ErasesEnv{}\ env\ bo\ \Gamma_{spec}\ t_0$ \\
--- (no analogue) & $\Lower{}\ \Gamma_{spec}\ t_0\ t$ and $\LowerEnv{}\ \Gamma_{spec}\ \Gamma$ \\
$\Sigma \vdash t \Rightarrow v$ (Rocq \code{eval}) & $\SEval{}\ env\ bo\ Us\ fl\ [\,]\ e\ v$ \\
$\Sigma' \vdash t' \Rightarrow v'$, squashed & $\WcbvEval{}\ \Gamma\ \code{eraseFlags}\ t\ v'$ \\
\code{firstorder\_ind} $\Sigma\ i$ & \code{FirstOrderInd} $env\ I$ (two clauses added) \\
\code{firstorder\_erases\_deterministic} (p.~8:63) & \code{firstorder\_erases\_deterministic} \\
\code{erase\_correct\_firstorder} (p.~8:63) & \code{shipping\_erase\_correct\_firstorder} \\
\code{axiom\_free} $\Sigma$ & no analogue in T5 \\
\end{tabular}
\end{center}

Five differences deserve to be stated plainly.

\emph{The target semantics is this repository's own \WcbvEval{}}, at \code{eraseFlags} --- prop-case
off, guarded fix, applied-form constructors --- documented as MetaRocq's \code{opt\_wcbv\_flags}. It is
a re-formalisation, not a transcription: Chapter~\ref{chap:lambdabox} records that the iota and
projection rules drop MetaRocq's parameter-count and saturation side conditions, so the relation is
strictly more permissive than MetaRocq's. A forward simulation into a more permissive target is a weaker
statement than one into the original.

\emph{The pass is part of T5's statement, but its own correctness is not what this chapter proves.}
MetaRocq's theorem relates one relation to one semantics; here the target term is the image of a
two-stage map, and \LowerEnv{} is an extra premise. The corresponding proof burden is the inversion kit
of Lemmas~\ref{lem:Lower-source_app}, \ref{lem:Lower-source_constApp} and \ref{lem:Lower-const_body}
and, above all, Proposition~\ref{prop:erases_elimSpine_no_value}. Conversely, T5 does not say the pass
preserves meaning in isolation; it says the \emph{composite} simulates the source. The pass's own
metatheory is Chapter~\ref{chap:lowering}.

\emph{Erasure is a relation and the theorem quantifies over its images.} MetaRocq has an erasure
function and states the relation as an extension of it. Here there is no function at this level, so
\code{Simulates} is universally quantified over $t_0$ and $t$ --- which is what makes the delta arm's
one-move treatment of the recursive and non-recursive cases possible, and what makes the first-order
collapse necessary rather than merely convenient.

\emph{The observation is first-order, and its domain is narrower than [S]'s.} \code{FirstOrderInd} adds
monomorphism and index-freedom to Letouzey's condition, so \code{List} and \code{Prod} fall outside it
even though \code{List Nat} and \code{Nat} $\times$ \code{Nat} are first-order data. And the checker
that decides the condition on \code{Nat}, \code{Bool} and a benchmark tree is sound only against the
\emph{table}, not against the model (Lemma~\ref{lem:firstOrderIndB_step}); the only proved model
inhabitant is an empty inductive (Lemma~\ref{lem:FOModel-firstOrderInd_E}).

\emph{No \code{axiom\_free} premise, and a different reason.} [S] restricts
\code{erase\_correct\_firstorder} to axiom-free environments because its proof runs through progress.
T5 needs no such premise: a body-less, non-constructor, non-type-name, non-eliminator constant simply
heads no \SEval{} value (Chapter~\ref{chap:source}), so a run reaching an axiom has no source derivation
at all.

Finally, what this chapter does \emph{not} establish. T5 is conditional on \code{UpstreamAsks\ env},
four lean4lean facts the pin does not prove; its measured axiom footprint contains \code{sorryAx},
inherited from the pinned lean4lean; and the published statement discards the environment accumulator
its own proof establishes. The observation of Section~\ref{sec:correctness:firstorder} is, at the level
of the model, inhabited only at an empty inductive. Chapter~\ref{chap:trust} collects these and the rest
of the trust boundary.
